%%%%%%%%%%%%%%%%%%%%%%%%%%%%%%%%%%%%%%%%%%%%%%%%%%%%%%%%%%%%%%%%%%%%%%
% LaTeX Example: Project Report
%
% Source: http://www.howtotex.com
%
% Feel free to distribute this example, but please keep the referral
% to howtotex.com
% Date: March 2011 
% 
%%%%%%%%%%%%%%%%%%%%%%%%%%%%%%%%%%%%%%%%%%%%%%%%%%%%%%%%%%%%%%%%%%%%%%

% Edit the title below to update the display in My Documents
%\title{Project Report}
%
%%% Preamble
\documentclass[paper=a4,12pt]{article}
\usepackage[T1]{fontenc}
\usepackage{kpfonts}
%\usepackage[fontsize=12pt]{scrextend}
\usepackage[utf8]{inputenc}
\usepackage[french]{babel}

% English language/hyphenation

\usepackage{amsmath,amsfonts,amsthm} % Math packages
%\usepackage[pdftex]{graphicx}	
\usepackage{url}
\usepackage[bottom=10em]{geometry}
\usepackage{float}
\usepackage{xcolor}
\usepackage{enumitem}
\usepackage{rotating}
\usepackage{lscape}
 


%%% Custom sectioning
%\usepackage{sectsty}
%\allsectionsfont{\normalfont\scshape}

%% Language definition package (for XML Annexe)
\usepackage{listings}
\usepackage{color}

%% Local modification of margins
%\newenvironment{changemargin}[2]{\begin{list}{}{%
%      \setlength{\topsep}{0pt}%
%      \setlength{\leftmargin}{0pt}%
%      \setlength{\rightmargin}{0pt}%
%      \setlength{\listparindent}{\parindent}%
%      \setlength{\itemindent}{\parindent}%
%      \setlength{\parsep}{0pt plus 1pt}%
%      \addtolength{\leftmargin}{#1}%
%      \addtolength{\rightmargin}{#2}%
%    }\item }{\end{list}}
%%%

%%% Custom headers/footers (fancyhdr package)
%\usepackage{fancyhdr}
%\pagestyle{fancyplain}
%\fancyhead{}											% No page header
%\fancyfoot[L]{}											% Empty 
%\fancyfoot[C]{}											% Empty
%\fancyfoot[C]{\thepage}									% Pagenumbering
%\renewcommand{\headrulewidth}{0pt}			% Remove header underlines
%\renewcommand{\footrulewidth}{0pt}				% Remove footer underlines
%\setlength{\headheight}{13.6pt}


%%% Equation and float numbering
\numberwithin{equation}{section}		% Equationnumbering: section.eq#
\numberwithin{figure}{section}			% Figurenumbering: section.fig#
\numberwithin{table}{section}				% Tablenumbering: section.tab#

%Graphics path
%\graphicspath{./Images/}

%%% Maketitle metadata
\newcommand{\horrule}[1]{\rule{\linewidth}{#1}} 	% Horizontal rule

\title{
  %\vspace{-1in} 			
  %\usefont{OT1}{bch}{b}{n}
  \horrule{1.5pt} \\[0.5cm]	
  \Huge \textbf{Projet Rain Of Music : \\ État de l'art} \\ [20pt]
 % \huge TM, \\ Université d'Osaka \\ [15pt]
 	%	\vspace{1cm}  
    \LARGE Année scolaire 2015-2016 \\ 
  \horrule{1.5pt} \\[0.5cm]
  %
}

\author{						
    \LARGE \underline{Encadrants} : Jean-Michaël Celerier, \\
   					\LARGE	\hspace{5cm} Myriam De Sainte-Catherine\\			
   	\vspace{1cm} 
   	\normalfont
   	\LARGE Élèves :  Akané LEVY, Maxime PAILLASSA    
}
\date{}

%%% Begin document
\begin{document}
\graphicspath{{./imgs/}{.}}
\maketitle

\begin{figure}[H]
  \centering\includegraphics[scale=1.2]{logo_enseirb.png}
\end{figure}

\newpage

\tableofcontents

\newpage
\normalsize
%\begin{changemargin}{-1cm}{-1cm}

\section*{Introduction}
- présenter les différents acteurs
- présenter le projet dans sa globalité (but final au delà du cadre de notre projet)
- expliquer brièvement ce que l'on va faire nous

\section{Objectifs}
Avant de nous intéresser plus en détail à ce que nous pouvons utiliser pour atteindre nos objectifs, nous allons préciser ceux-ci.

Comme dit précédemment, notre travail consistera à écrire un logiciel permettant de visualiser et de simuler une chorégraphie écrite par les artistes. Cette chorégraphie est écrite grâce au logiciel i-score, qui est un séquenceur: il permet d'écrire des scénarios en programmant des évènements dans le temps. Notre travail impliquera donc deux problématiques: 
\begin{itemize}
\item la récupération des données d'i-score dans le logiciel, de manière à avoir les positions des robots.
\item l'affichage des robots dans une scène 3D, qui impliquera d'autres problématiques comme la gestion de collisions.
\end{itemize} 




\section{Existant}
\subsection{i-score}

Tout le travail réalisé devra s'interfacer avec le logiciel i-score qui sera utilisé pour écrire les chorégraphies. Pour communiquer avec un autre logiciel, i-score intègre le format OSC qui utilise UDP. Un message OSC se compose comme ceci:
\begin{lstlisting}
<address pattern> <data type> <data>
\end{lstlisting}
Address pattern est l'URL désignant la destination du message, et data type une chaîne de caractères indiquant le type de données contenues dans data.
Cependant, le protocole Minuit peut aussi être utilisé: ce protocole se situe juste au-dessus du protocole OSC, et permet de mieux organiser les données transmises. Les adresses OSC sont organisées dans un arbre et le logiciel pourra demander via le protocole Minuit à i-score de lui communiquer la valeur de tel ou tel paramètre.  


\subsection{Visualisation et simulation}
Cette partie présentera quelques articles scientifiques se rapprochant de notre projet autour des thématiques telles que la visualisation de scène 3D d'une peformance scénique et la modélisation de robot dans un espace 3D. 

Dans le domaine de visualisation de performance scénique 3D, un des objectifs est d'obtenir un rendu se rapprochant le plus possible de la réalité de façon à ce que l'on puisse avoir un aperçu réaliste de la performance conçue virtuellement. C'est dans cette optique que le logiciel de visualisation et simulation 3D StageViz\cite{StageViz} a été développé . Ce logiciel permet de visualiser en 3D en temps réel une performance scénique qui a été conçue virtuellement à l'aide d'un autre logiciel dédiée à cela. Dans la forme, nous retrouvons ici, un cas très similaire à notre objectif avec la présence de visualisation mais aussi d'interopérabilité avec un outil de conception scénique.

StageViz se base sur trois éléments principaux que l'on doit aussi retrouver dans notre projet : les objets 3D, la manipulation de la scène et la chronologie. Dans le cadre de notre projet, les objets 3D représenteront les robots et la scène pourra être visionner sous tous les angles grâce à une camera mobile. En ce qui concerne la composante temporelle, elle sera gérée au niveau d'i-score qui nous envoie les données sur les robots au fur et à mesure. 

Une grande différence entre StageViz et notre projet réside dans l'utilisation du rendu 3D. StageViz serait un outil pour une simulation réaliste pour visionner les différents états des objets dans la scène au cours de la performance. Notre but serait de visualiser l'évolution des positions des objets dans la scène pour aider l'écriture de la chorégraphie. L'idée est donc d'avoir un rendu permettant une vision simplifiée et globale de la chorégraphie avec tous les robots.

Dans le domaine de performance par les robots, une équipe de recherche a étudié la modélisation des différents mouvements d'un robot dans le but de visualiser une performance \cite{robotArt}. Le robot n'est pas considéré en tant qu'un seul objet avec une position mais comme un ensemble de composants possédant des paramètres (degré de liberté) et organisés hiérarchiquement. Ainsi lorsqu'un mouvement est exécuté au niveau d'une composante, il est appliqué à toutes ses composantes filles.
%Il existe déjà des logiciels permettant de réaliser un résultat semblable à ce qui est attendu. Il y a notamment le logiciel StageViz pour la visualisation en 3D en temps réel une performance scénique, qui peut être utilisé avec des outils de conception de performance scénique. Nous retrouvons ici, un cas très similaire à notre objectif avec la présence de visualisation mais aussi d'interopérabilité avec un outil de conception scénique.

Les auteurs de StageViz expliquent qu'à la base, un sondage avait été réalisé auprès des différentes personnes du milieu de la performance artistique (artistes, ingénieurs, metteur en scène) pour cibler les besoins d'un outil de visualisation. Trois éléments principaux ont été retenus : les objets, la manipulation de la scène et la chronologie. En effet, la présence des différents objets 3D, les décors mais aussi les acteurs, permettent un rendu réaliste de la scène. Concernant la manipulation de la scène, c'est ce qui représente l'avantage d'avoir une modélisation virtuelle : la possibilité de bouger la camera dans toute la scène 3D pour avoir un aperçu sous tous les angles. La chronologie permet d'avoir une évoution des états des objets 3D dans la scène.

Contrairement à StageViz dont le but semble être de simuler de façon la plus réaliste possible une performance scénique, le but de notre projet repose plus sur une visualisation pour aider l'écriture de la chorégraphie. L'idée est donc d'avoir un rendu réaliste mais surtout permettre une vision simplifiée et globale de la chorégraphie avec tous les robots.

%- traite sur les performances, 
%	parle de "storyboarding" -> notion de temps 
%	
%- visualisation de la performance avec des composantes créées par l'utilisateur
%BUT LOGICIEL "StageViz provides visualized scene image
%	based on the timeline and virtual scene components created by
%	a user." mais dans les articles précédents, expliqué que le logociel permet aussi de contrôler différents outils liés à la mise en scène (SFX, lights...)
%
%- BUT ARTICLE :  établir une connection entre la visualisation et la performance réelle sur scène
%
%- pourquoi ce logiciel -> sondage auprès des gens du milieu de la performance artistique (ingénieurs, artistes, metteur en scène)
%	trois critères ont été relevés et retenus
%		objet / "timeline"~chronologie / manipulation de la scène 
%
%- logiciel basé sur Direct3D pour faire le rendu 3D
%
%- expérience utilisant des costumes lumineux 
%	  résultat : le rendu 3D permet de faire une modélisation très fine de la réalité 
%		  montre que ce logiciel peut être utilisé pour créer/réaliser une mise en scène puisque le rendu visuel généré par ce logiciel est très proche de ce qui peut être obtenu en vrai.
%Dans cet article, les auteurs proposent un système d'assemblage de robot dans un environnement virtuel en 3D. Le but est ensuite de pouvoir les mettre en mouvement. Chaque robot est composé de différentes parties organisées hiérarchiquement. Pour faire bouger les différentes parties des robots assemblés, le module de visualisation 3D proposé se base sur des matrices de transformation (translation, rotation, facteur d'échelle). Quand une partie subit une transformation, celle-ci est appliquée à chacune des parties filles. \\
Une interface permet à l'utilisateur d'assembler les robots et de les observer depuis quatre points de vue différents. 




\section{Bibliothèques 3D}
Dans cette partie, nous allons présenter différentes bibliothèques graphiques 3D que nous pourrions utiliser: OpenSceneGraph, Ogre, OpenFrameworks, irrlicht, Blender, Unity, Qt3D, Babylon.js et Three.js. \\
Toutes ces biliothèques exceptées les cinq dernières utilisent OpenGL et sont déployables sur pratiquement toutes les plateformes (y compris smartphones et tablettes). Le fait qu'elles soient orientées objets permet d'avoir une programmation plus haut niveau tout en profitant des performances d'OpenGL. Nous n'avons pas voulu utiliser uniquement OpenGL car cela aurait demandé plus de code et de temps pour un résultat équivalent. Elles intègrent également des packages additionnels pour faire des rendus 3D plus poussés. \\
Blender et Unity sont plus orientées sur de l'animation et du jeu vidéo donc sur du rendu 3D. \\
Babylon.js et Three.js se déploient sur navigateur via javascript et utilisent WebGL.

- vérifier la version de C++ utilisée \\
- communautés dév (bibliothèque maintenue et contribution des utilisateurs) et utilisateurs (domaines d'applications variés) \\

Pour chaque bib : 
- particularité / explication brève
- domaine d'utilisation
- moteur physique 
- philosophie
- license
- communauté dév et utilisateurs
- stabilité
- desavantages ?

\subsection{OpenSceneGraph}
- explication : open source/C++/
philosophie/choix :  manière différente de visualisation -> difficulté à prendre en main (désavantage pour nous et pour les suivants) > difficulté de maintenabilité
- communauté : peu active

\subsection{OpenFrameworks}
- explication : open source/C++/

%- fonctionne avec OpenGL, OpenCV (traitement image pour trajectoires?)
- philosophie : "it should be collaborative, usable and simple, consistent and intuitive, cross-platform, powerful, and extensible. openFrameworks is also driven by a "do it with others" (DIWO) philosophy"
- désavantage installation difficile ?


\subsection{Ogre}
- explication : open source/C++/
domaine d'utilisation :  plus tourné jeux vidéos 
license : GNU LGPL \\

\subsection{Irrlicht}
- explication : open source/C++/
moteur physique : oui avec détection de collision
stabilité : pas très stable 

\subsection{Blender et Unity}
- explication : Blender : open source
				Unity : C# / pas open source
				
- communauté  : utilisateur active
- domaine d'utilisation : tourné jeux vidéos mais surtout modélisation et pas vraiment visualisation
--> peut être trop de fonctionnalités liées à la modélisation qui nous seraient inutiles...?
- désavantage pour Unity -> pas open source

\subsection{Qt3D}
- explication : open source/ C++
moteur physique : détection de collision
stabilité : pas vraiment stable = en cours de développement

\subsection{Babylon.js et Three.js}
- explication: moteur 3D javascript utilisant webgl, permet d'avoir un rendu visuel dans le navigateur
- désavantage : 
		besoin de mettre dans navigateur  ? 
	javascript : pas adapté pour la suite du projet, a priori le projet dure plus que nos trois mois (environ un an ou deux) donc si on considère le fait que la visualisation devienne sophistiquée (par ex : possibilité de modifier la chorégraphie en elle-même ? pas seulement visualisation mais aussi création) le langage javascript ne semble pas adapté...
		GROS NEGATIF : difficulté à récupérer les données d'iscore via UDP avec Javascript? (utiliser node.js ?)
%Dans cette partie, nous allons présenter les bibliothèques graphiques 3D que nous pourrions utiliser: OpenSceneGraph, Ogre, OpenFrameworks, irrlicht \\
Toutes ces biliothèques utilisent OpenGL et sont déployables sur pratiquement toutes les plateformes (y compris smartphones et tablettes). Le fait qu'elles soient orientées objets permet d'avoir une programmation plus haut niveau tout en profitant des performances d'openGL. Elles intègrent également des packages additionnels pour faire des rendus 3D plus poussés. \\

- vérifier la version de C++ utilisée \\
- communautés dév (bibliothèque maintenue et contribution des utilisateurs) et utilisateurs (domaines d'applications variés) \\

- Ogre, Osg et irrlicht: graphes \\


irrlicht: + détection de collision, - pas très stable \\
ogre : - plus tourné jeux vidéos 

- philosophie: \\
OF:  it should be collaborative, usable and simple, consistent and intuitive, cross-platform, powerful, and extensible. openFrameworks is also driven by a "do it with others" (DIWO) philosophy. \\




\section{Bilan et choix de l'outil}
liste des critères:
\begin{itemize}
\item base 3D (gestion caméra, scène)
\item open source license GPL
\item détection de collisions (moteur physique)
\item auto-suffisant (facilité d'utilisation d'UDP)
\item unicité du langage avec i-score (C++)
\item multi-plateforme
\item taille adaptée au projet (pas de surplus de fonctionnalités)
\item maintenabilité (bonne communauté active)
\item maintenabilité 2 (facilité à prendre en main)
\end{itemize}
%justification

\section*{Conclusion}
- re-justification de la bibliothèque choisie en reprenant le bilan précédent
- expliquer commment on compte utiliser l'outil choisi pour notre projet


  \newpage

  \bibliographystyle{plain}
  \bibliography{bibli.bib}
  \newpage

  \appendix
  \part*{Annexes}
  \input{annexe.tex}
 %   \section{Autres}% \label{annexe:modeles}
    %\input{notice_installation.tex} 
  

  
 % \newpage 
  

 

%\end{changemargin}

%%% End document
\end{document}
